\section{Why an equilibrium model alone has no mixing time}\label{sec:specification}
The results in the main text specify both the update rule and the
clock. Two elementary constructions show why these are mathematical
assumptions, even if the equilibrium law and reversibility are already
fixed. A third calculation distinguishes high-probability asymptotic
bounds from essential bounds over every disorder realization.

\begin{proposition}[Arbitrary refresh times]\label{prop:arbitrary-refresh}
Fix any realization of the SK measure on $\Omega_n$, $n\ge2$. For every positive integer $r$, there is an irreducible, aperiodic, reversible Markov kernel with stationary law $\mu$ whose worst-state threshold-$1/4$ mixing time equals $r$.
\end{proposition}
\begin{proof}
At each step, refresh the whole configuration from $\mu$ with
probability $p$ and otherwise leave it fixed. In operator notation,
$\Pi(x,y)=\mu(y)$ and $K_p=(1-p)\mathrm{Id}+p\Pi$, where $0<p\le1$.
All entries of $K_p$ are positive, and
$\mu(x)K_p(x,y)=\mu(y)K_p(y,x)$. The identity $\Pi^2=\Pi$ gives,
for $0<p<1$,
\[
 K_p^k=\Pi+(1-p)^k(\mathrm{Id}-\Pi),\qquad
 \max_x\|K_p^k(x,\cdot)-\mu\|_{\TV}
 =(1-p)^k(1-\mu_{\min}).
\]
Here $\mu_{\min}\le2^{-n}\le1/4$. Choosing
\[
 1-p=\left[\frac{1}{4(1-\mu_{\min})}\right]^{1/r}
\]
gives distance exactly $1/4$ at $r$ and strictly larger at every earlier integer. The case $p=1$ also has mixing time one.
\end{proof}

The same freedom remains if transitions are required to change at
most one spin. Given the uniform-site heat-bath kernel $P$, the kernel $(1-p)\mathrm{Id}+pP$ has the same stationary law, detailed balance, irreducibility, and transition restriction. Starting from a minimum-mass state, its distance after $k$ steps is at least $(1-p)^k-\mu_{\min}$, by the event that no update was attempted. Its mixing time consequently diverges as $p\downarrow0$. Allowing $p$ to depend on $n$ changes possible growth classes as well.

We next fix the dynamics completely. The remaining distinction is
between typical disorder as $n$ grows and all disorder at a fixed
dimension. Unbounded Gaussian couplings already separate those notions
for two spins.

\begin{proposition}[A two-spin example]\label{prop:two-spin}
For $n=2$ and the discrete uniform-site heat-bath kernel, the mixing time has no finite deterministic essential upper bound over the Gaussian coupling, although it is finite for every realization.
\end{proposition}
\begin{proof}
A large positive coupling keeps the two aligned states persistent.
This can be quantified using the linear observable $g(x)=x_1+x_2$.
Put $w=W_{12}$. Spin-flip symmetry gives $\mu[g]=0$, and conditional
expectation computes the eigenvalue exactly:
\[
 Pg=\theta g,\qquad \theta=\frac{1+\tanh w}{2}.
\]
Starting at $(1,1)$, the expectation of $g$ after $k$ steps is $2\theta^k$. Since $\|g\|_\infty=2$, the dual total-variation bound yields
\[
 \|P^k((1,1),\cdot)-\mu\|_{\TV}\ge\tfrac12\theta^k.
\]
For each fixed $k$, this exceeds $1/4$ when $w$ is sufficiently large and positive. Every such Gaussian tail has positive probability. The chain for any finite $w$ remains finite, irreducible and aperiodic, so its own mixing time is finite.
\end{proof}

Proposition~\ref{prop:two-spin} concerns essential bounds over all disorder realizations at fixed $n$. It does not contradict the typical-disorder asymptotic statements in Theorem~\ref{thm:mixing}.
